\begin{table}[!tbp]\centering
\caption{Robustez frente a la ola migratoria venezolana}\label{tab:migration}
\begin{threeparttable}
\footnotesize
\setlength{\tabcolsep}{4pt}
\begin{tabular}{lcccc}
\toprule
Especificación & Log salario/hora & Empleo & Empleo formal & Trab. independiente \\
\midrule
Base & -0.021 & -0.025$^{***}$ & -0.046$^{***}$ & 0.036$^{***}$ \\
 & (0.012) & (0.006) & (0.007) & (0.006) \\
EF departamento$\times$trimestre & -0.020 & -0.015$^{**}$ & -0.041$^{***}$ & 0.034$^{***}$ \\
 & (0.013) & (0.006) & (0.008) & (0.006) \\
Excluyendo Lima y Callao & -0.021 & -0.018$^{*}$ & -0.042$^{***}$ & 0.037$^{***}$ \\
 & (0.022) & (0.009) & (0.010) & (0.010) \\
Sin los 8 deptos. con más migrantes & -0.051 & -0.012 & -0.049$^{***}$ & 0.045$^{**}$ \\
 & (0.032) & (0.013) & (0.015) & (0.016) \\
\bottomrule
\end{tabular}
\begin{tablenotes}\footnotesize
\item Notas: Estimaciones DiD de $\mathrm{Low}\times\mathrm{Post}$ bajo especificaciones que abordan la llegada de inmigrantes venezolanos. Los efectos fijos departamento$\times$trimestre absorben cualquier choque local común a ambos grupos educativos dentro de un departamento-trimestre, incluidas las llegadas de migrantes y las olas de regularización. La exclusión de los 8 principales elimina Lima, Callao, La Libertad, Arequipa, Lambayeque, Piura, Ica y Tumbes, los departamentos que acogen a la gran mayoría de la población venezolana. Errores estándar agrupados por departamento; factores de expansión de la ENAHO. $^{*}p<0.1$, $^{**}p<0.05$, $^{***}p<0.01$.
\end{tablenotes}
\end{threeparttable}\end{table}
